\documentclass{article}
\usepackage{amsmath, amssymb, amsthm}
\usepackage{hyperref}
\usepackage{geometry}
\geometry{margin=1in}

\title{Erd\H{o}s Problem \#376}
\date{Last updated: 2025-08-31}
\author{Source: erdosproblems.com}

\begin{document}
\maketitle

\noindent\textbf{Status:} OPEN \quad \textbf{No prize}

\noindent\textbf{Tags:} number theory, binomial coefficients, base representations

\medskip

\noindent\textbf{Problem Statement:}

Are there infinitely many $n$ such that $\binom{2n}{n}$ is coprime to $105$?




Erd\H{o}s, Graham, Ruzsa, and Straus \cite{EGRS75} have shown that, for any two odd primes $p$ and $q$, there are infinitely many $n$ such that $\binom{2n}{n}$ is coprime to $pq$. This is equivalent (via Kummer's theorem ) to whether there are infinitely many $n$ which have only digits $0,1$ in base $3$, digits $0,1,2$ in base $5$, and digits $0,1,2,3$ in base $7$. The sequence of such $n$ is A030979 in the OEIS. The best result in this direction is due to Bloom and Croot \cite{BlCr25}, who proved that, if $p_1,p_2,p_3$ are sufficiently large primes, then there are infinitely many $n$ such that almost all of the base $p_i$ digits are $&lt;p_i/2$. In other words, for all $\epsilon&gt;0$, there are infinitely many $n$ such that $\binom{2n}{n}$ is coprime to $p_1p_2p_3$, except for a factor of size $\leq n^\epsilon$. This is mentioned in problem B33 of Guy's collection \cite{Gu04}. It is also discussed in an article of Pomerance \cite{Po15c}. Graham offered \$1000 for a solution to this problem (as mentioned in \cite{Gu04} and \cite{BeHa98}).




References

[BeHa98] Berend, Daniel and Harmse, J\o rgen E., On some arithmetical properties of middle binomial coefficients . Acta Arith. (1998), 31--41.

[BlCr25] T. F. Bloom and E. Croot, Integers with small digits in multiple bases . arXiv:2509.02835 (2025).

[EGRS75] Erd\H{o}s, P. and Graham, R. L. and Ruzsa, I. Z. and Straus, E. G., On the prime factors of $(\sp{2n}\sb{n})$ . Math. Comp. (1975), 83-92.

[Gu04] Guy, Richard K., Unsolved problems in number theory . (2004), xviii+437.

[Po15c] Pomerance, Carl, Divisors of the middle binomial coefficient . Amer. Math. Monthly (2015), 636--644.

\noindent\textbf{Related OEIS sequences:} A030979


\bigskip
\noindent\small{Source: \url{https://www.erdosproblems.com/376}}
\end{document}
